\begin{abstract}
\begin{itemize}
    \item In recent year, the power of neural networks for solving combinatorial problems has been investigated. 
    \item Many interesting combinatorial problems are NP-hard, and therefore they do not admit efficient algorithms under the assumption that is P is not equal to NP. 
    \item It could still be the case that graph neural networks may solve instances of practical relevance. 
    \item The clique-colouring problem consists of determining whether a graph has a k-clique or whether the graph is (k-1)-colorable. 
    \item We show that graph neural networks are Unable to disguighs graphs that have a k-clique from graphs that are (k-1)-colorable. 
    \item This is an unconditional result. We exhibit an infinite family of graphs that cannot be distinguished by any graph neural network. 
    \item We generalise this result to higher-dimensional graph neural networks. 

\end{itemize}
\end{abstract}